% =======================================================================================
% Chapter 2 -- Literature Review
% SECTION-H(11): theoretical background, previous work, gaps in knowledge.
% =======================================================================================
\chapter{Literature Review}
\label{ch:literature}

\section{Cascaded and Direct Speech-to-Speech Translation}
\label{sec:lit_s2st}

Speech-to-speech translation has been approached in two ways. The \emph{cascaded} approach chains
automatic speech recognition, machine translation and speech synthesis, so that each stage can be
built, replaced and measured independently. The \emph{direct} approach trains a single model to map
source speech to target speech, either through an intermediate spectrogram
\cite{jia2022translatotron2} or through discrete acoustic units learned without text
\cite{lee2022discreteunits,inaguma2023unity}.

The trade-off is well documented. Direct models avoid the propagation of recognition errors into
translation and can, in principle, carry paralinguistic information that text discards; cascades
remain competitive in translation quality, are far cheaper to assemble, and are interpretable at
every stage. \citeA{sperber2020endtoend} argue that most of the theoretical advantages claimed for
end-to-end systems have yet to be realised in practice, and \citeA{bentivogli2021cascade} find the
remaining quality difference on standard benchmarks to be small. For a study whose object is
\emph{measurement} rather than architecture, the cascade has a decisive advantage: the intermediate
text makes it possible to attribute an observed error to a specific stage, and to quantify what
each stage costs. That is the argument for the design adopted here.

\section{Automatic Speech Recognition}
\label{sec:lit_asr}

\citeA{radford2023whisper} trained the Whisper family on approximately 680{,}000 hours of weakly
supervised multilingual audio, showing that scale alone yields zero-shot robustness competitive
with supervised systems. Whisper also emits segment-level timestamps, which matters here: dubbing
continuous video requires knowing where each utterance sits on the timeline, not merely what was
said. Its known weaknesses --- disfluency, overlapping talk, short backchannels --- are exactly the
conditions of the conversational corpus used here, and the resulting error rate is reported rather
than suppressed.

\section{Neural Machine Translation and Length Control}
\label{sec:lit_mt}

Transformer-based translation \cite{vaswani2017attention} is standard. The Marian toolkit
\cite{junczysdowmunt2018marian} and the OPUS-MT release \cite{tiedemann2020opusmt} provide small,
fast, openly licensed bilingual models trained on OPUS --- books, subtitles and institutional
documents. Their register is written text, a documented mismatch for spontaneous dialogue;
in-domain fine-tuning is the standard remedy and is examined empirically in
Chapter~\ref{ch:results}.

For dubbing specifically, translation quality is not the only requirement: the translated utterance
must also be \emph{speakable in the time available}. \citeA{federico2020dubbing} decompose automatic
dubbing into translation, prosodic alignment and audio rendering, and \citeA{lakew2021verbosity}
show that verbosity can be controlled at translation time by conditioning on a desired length
class. This literature anticipates the duration problem reported in
Section~\ref{sec:duration_result} and points to its proper solution.

\section{Speaker-Identity-Preserving Speech Synthesis}
\label{sec:lit_tts}

Neural text-to-speech moved from Tacotron~2 with a neural vocoder
\cite{shen2018tacotron2,kong2020hifigan} to fully end-to-end variational models
\cite{kim2021vits}. Multi-speaker synthesis was made zero-shot by conditioning the synthesiser on
an embedding produced by a separately trained speaker-verification network
\cite{jia2018sv2tts}: because the encoder generalises to unseen speakers, a few seconds of reference
audio is enough to clone a voice. YourTTS \cite{casanova2022yourtts} extended this to the
multilingual case, and XTTS-v2 \cite{casanova2024xtts} scales the idea to seventeen languages with
short reference clips. Voice conversion research reaches similar capability from another direction,
by disentangling content from speaker information in an autoencoder \cite{qian2019autovc}.

Two properties of this family matter for the present study. First, conditioning is on
\emph{timbre}: the speaker embedding is trained for verification, and it deliberately discards
utterance-level variation. Second, nothing in the architecture supplies the source utterance's pitch
contour to the decoder. The synthesiser must therefore infer intonation from the target text alone.
This is the mechanism behind the result reported in Section~\ref{sec:prosody_gap}, and it is a
property of the design rather than a defect of the implementation.

\section{Prosody Modelling and Cross-Lingual Prosody Transfer}
\label{sec:lit_prosody}

Within a single language, prosody transfer is comparatively mature.
\citeA{skerryryan2018prosody} learn a latent prosody embedding from a reference recording and show
that conditioning Tacotron on it reproduces the reference's prosody with fine time detail, even
across speakers. Global Style Tokens \cite{wang2018gst} obtain interpretable style control without
prosody labels, and \citeA{yang2025emotional} surveys emotional synthesis and voice conversion in
the same tradition. Reviews of multilingual voice-to-voice translation
\cite{ambadekar2025review} and of multi-feature multilingual synthesis
\cite{nuthakki2023multilingual} converge on the same conclusion: translation, synthesis and
speaker-preserving components must be integrated, and prosody is the component most often lost.

Cross-lingual prosody transfer is markedly harder, and the obstacle is data. Transferring prosody
from a source-language recording to target-language speech requires supervision about what the
target-language rendition should sound like, and parallel expressive recordings of the same content
in two languages are rare. \citeA{song2023styles2st} address this with zero-shot style transfer
inside a direct S2ST model. \citeA{peng2024mslm} preserve speaker style in a textless speech
language model. The Seamless system \cite{seamless2023} adds explicit expressivity and streaming to
large-scale multilingual translation. Closest to the present work,
\citeA{swiatkowski2023dubbing} build a prosody representation for machine dubbing that is
transferable across both language and speaker, using a variational reference encoder with a noise
modelling module for field recordings --- a design necessitated precisely by the absence of clean
parallel expressive data.

The corpus used in this dissertation removes that obstacle for measurement purposes.
\citeA{ward2023dral} designed DRAL to support ``machine learning of cross-language prosodic
mappings'': participants re-enact their own utterances in the other language, so both sides of a
pair come from the same speaker performing the same communicative act. What is a scarcity problem
for \emph{training} an expressive dubbing model is not a scarcity problem for \emph{evaluating} one.

\section{Evaluation of Cross-Lingual Speech Generation}
\label{sec:lit_eval}

Translation quality is measured with BLEU \cite{papineni2002bleu} and chrF
\cite{popovic2015chrf} through a standardised implementation \cite{post2018sacrebleu}, and
recognition with word error rate. Speaker similarity is a cosine between speaker-verification
embeddings; ECAPA-TDNN \cite{desplanques2020ecapa}, trained on VoxCeleb \cite{nagrani2017voxceleb}
and distributed through SpeechBrain \cite{ravanelli2021speechbrain}, is the current standard.
Prosody is described with fundamental frequency, energy, duration and speaking rate, F0 being
estimated by pitch trackers such as pYIN \cite{mauch2014pyin} in librosa \cite{mcfee2015librosa}.

Three methodological weaknesses recur and are addressed explicitly in
Chapter~\ref{ch:methodology}. Speaker-similarity cosines are reported without a same-speaker
cross-language ceiling, so their scale is undefined. Intrusive signal metrics such as PESQ and STOI
are proposed although they require a time-aligned reference of \emph{the same utterance}, which
cannot exist when the output is in another language. And correlations are reported pooled across
speakers, counting stable between-speaker differences as evidence of transfer; reporting the
within-speaker component alongside is the standard correction, adopted here.

\section{Research Gap}
\label{sec:gap}

The literature establishes that zero-shot cloning preserves timbre, that within-language prosody
transfer is achievable, and that cross-lingual prosody transfer is limited chiefly by parallel
expressive data. What it does not establish is a calibrated statement of \emph{how much} identity
and \emph{how much} prosody a standard, openly available cascade actually carries when the human
answer for the same speaker and the same content is available for comparison.

This dissertation fills that gap. It uses a same-speaker parallel corpus as a measuring instrument:
identity is expressed as a fraction of the human cross-language ceiling; prosody transfer is
expressed against the human within-speaker correspondence for the same features; and the
predictable component of cross-lingual prosody is isolated by testing a supervised mapping against
baselines that a useful model must beat. The contribution is therefore methodological and
empirical rather than architectural, and its value is that it states, in reusable units, the size
of the gap that expressive dubbing systems set out to close.
